% Purpose: Appendix F, sensitivities of the network: runs that reach the epoch limit (Table F1) and the direction of
% their effect on A, the high-moment exclusion and the untuned CGCNN check. Converted from the Supplementary
% Materials, Sections S2.3 and S6 (paper/si/sections/s2_training.tex, s6_sensitivities.tex), in session A4, with the
% final-network records only. Figure F1 (fig:app_sens; paper/figgen/figureF1_sensitivities.py). Input by main.tex after \appendix. Numbers:
% results/derived_w2/w2_training_limit.json (docs/w2_diagnostics.md (b)), results/precision_a_b400_b654.json,
% results/derived_n4/high_moment_final_network.json, results/sens_high_moment_rf.json, results/sens_cgcnn.json.
\section{Sensitivities}
\label{app:sens}

\subsection{Runs That Reach the Epoch Limit}

A run reaches the limit of \numNetCapMaxEpochs{} epochs (Appendix~\ref{app:network}) when early stopping has not ended it by the last epoch. In the learning-curve design this happens rarely for S and D (Table~\ref{tab:app_capped}); the descriptor predictor P1 has its own limit (Appendix~\ref{app:network}). Over the additional resamples with \numCfgBudgetsMax{} training hosts, \numPrecisionABFourHundredBSixHundredFiftyFourDescriptiveCappedPrecisionSixHundredFiftyFourVTwoSNCapped{} of \numPrecisionABFourHundredBSixHundredFiftyFourDescriptiveCappedPrecisionSixHundredFiftyFourVTwoSNRuns{} runs of S and \numPrecisionABFourHundredBSixHundredFiftyFourDescriptiveCappedPrecisionSixHundredFiftyFourVTwoDNCapped{} of \numPrecisionABFourHundredBSixHundredFiftyFourDescriptiveCappedPrecisionSixHundredFiftyFourVTwoDNRuns{} of D reach the limit.

\begin{table}[H]
\caption{Runs of the learning-curve design that stop at the limit of \numNetCapMaxEpochs{} epochs rather than by early stopping, out of \numNetCapRunsSixHundredFiftyFour{} runs (\numCfgSplitNOuterResamples{} resamples, three seeds) per training size. Record: \texttt{results/derived\_w2/w2\_training\_limit.json}.\label{tab:app_capped}}
\centering
\small
\begin{tabular}{rrr}
\toprule
Training hosts & S & D \\
\midrule
\numCfgBudgetsHostsItemOne{} & \numNetCapSTwentyFive{} & \numNetCapDTwentyFive{} \\
\numCfgBudgetsHostsItemTwo{} & \numNetCapSFifty{} & \numNetCapDFifty{} \\
\numCfgBudgetsHostsItemThree{} & \numNetCapSOneHundred{} & \numNetCapDOneHundred{} \\
\numCfgBudgetsHostsItemFour{} & \numNetCapSTwoHundred{} & \numNetCapDTwoHundred{} \\
\numCfgBudgetsHostsItemFive{} & \numNetCapSFourHundred{} & \numNetCapDFourHundred{} \\
\numCfgBudgetsHostsItemSix{} & \numNetCapSSixHundredFiftyFour{} & \numNetCapDSixHundredFiftyFour{} \\
\bottomrule
\end{tabular}
\end{table}

The runs that reach the limit lower $A$ slightly rather than raise it. With \numCfgBudgetsMax{} training hosts, the single-run site MAE of the runs of S that reach the limit minus that of the other runs of S is \numNetCapThirtyByFiveCappedMinusUncappedMaeS{}~eV over the additional resamples and \numNetCapSweepCappedMinusUncappedMaeS{}~eV in the learning-curve design: they are the runs still improving when training stopped. Pairs of S and D runs that contain such a run have a smaller $A$, so including them changes the mean single-run $A$ by \numNetCapThirtyByFiveEffectOnA{}~eV over the additional resamples and by \numNetCapSweepEffectOnA{}~eV in the learning-curve design, and removing them from the seed ensembles moves $A$ over the additional resamples from \numNetCapEnsembleAAll{} to \numNetCapEnsembleACappedRemoved{}~eV. Training S longer would therefore more likely shrink $A$ than grow it; the effect is small next to the interval of $A$ and is descriptive. Of the \numNetDclassRuns{} runs that read one descriptor class (Section~\ref{sec:disc_principal}), \numNetDclassCapRuns{} reach the limit; removing their \numNetDclassCapCells{} cells from every model leaves the gaps to D with \numCfgBudgetsMax{} training hosts at \numNetDclassSixHundredFiftyFourCapGapDHost{} and \numNetDclassSixHundredFiftyFourCapGapDSite{}~eV, against \numNetDclassSixHundredFiftyFourGapDHost{} and \numNetDclassSixHundredFiftyFourGapDSite{}~eV with all runs.

\subsection{Entries with a High Defect Moment}

The hosts are classified as non-magnetic in the release, but \numSensHighMomentRfNExcludedEntries{} neutral vacancy calculations of the dataset end with a defect magnetic moment above \numSensHighMomentVOneThresholdMuB{}~$\mu_\mathrm{B}$. All lie in La- or Ce-containing oxides, apart from one tungstate, that is, in hosts with a reducible cation on which the electrons left by the vacancy can localise instead of occupying a delocalised vacancy state. Such defects lie outside the deep, closed-shell vacancy that the rest of the dataset describes, so they test whether the conclusion depends on them. The check removes them from training and test and repeats the fit. For the network with \numCfgBudgetsMax{} training hosts (ten resamples, three seeds), compared on the same test sites, the removal moves $A$ from \numNetHighMomentCommonA{} [\numNetHighMomentCommonALo{}, \numNetHighMomentCommonAHi{}]~eV to \numNetHighMomentExclA{} [\numNetHighMomentExclALo{}, \numNetHighMomentExclAHi{}]~eV, a change of \numNetHighMomentChangeA{} [\numNetHighMomentChangeALo{}, \numNetHighMomentChangeAHi{}]~eV whose interval contains zero (Figure~\ref{fig:app_sens}a; \texttt{results/derived\_n4/high\_moment\_final\_network.json}). The one-sided upper bound on $A$ becomes \numNetHighMomentExclAUpper{}~eV, below $\delta = \numNetTolerance{}$~eV, so the statement at \numCfgBudgetsMax{} training hosts holds without these entries. The margin is narrow, and the two-sided interval no longer lies inside $\pm\delta$. For RF-Kumagai with \numSensHighMomentRfPoolSize{} training hosts the removal changes the site MAE by \numSensHighMomentRfDeltaTotalExcludedMinusBaselineMaeMean{} [\numSensHighMomentRfDeltaTotalExcludedMinusBaselineMaeCiLo{}, \numSensHighMomentRfDeltaTotalExcludedMinusBaselineMaeCiHi{}]~eV (\texttt{results/sens\_high\_moment\_rf.json}), because the entries make up a small part of every test set.

\begin{figure}[H]
\centering
\figph{figureF1_sensitivities}
\caption{Advantage $A$ without high-moment entries, and an untuned CGCNN beside the network. (\textbf{a}) $A$ of the network with \numCfgBudgetsMax{} training hosts on all test sites, on the common test sites and on the common sites without the entries, with two-sided intervals, one-sided upper bounds and the margin $\delta$. (\textbf{b}) Test MAE of S and D of the network and of CGCNN on one resample.\label{fig:app_sens}}
\end{figure}

\subsection{A Second Architecture}

An untuned CGCNN \citep{Xie2018crystal}, with fixed optimiser settings, was trained as S and D on resample \numSensCgcnnResample{} at \numSensCgcnnBudgetsItemOne{}, \numSensCgcnnBudgetsItemTwo{} and \numSensCgcnnBudgetsItemThree{} training hosts (\texttt{results/sens\_cgcnn.json}). Its advantage falls with the number of training hosts as that of the network does, from \numSensCgcnnTableCgcnnFiftyAPoint{} [\numSensCgcnnTableCgcnnFiftyACiNinetyFiveItemOne{}, \numSensCgcnnTableCgcnnFiftyACiNinetyFiveItemTwo{}]~eV at \numSensCgcnnBudgetsItemOne{} hosts to \numSensCgcnnTableCgcnnSixHundredFiftyFourAPoint{} [\numSensCgcnnTableCgcnnSixHundredFiftyFourACiNinetyFiveItemOne{}, \numSensCgcnnTableCgcnnSixHundredFiftyFourACiNinetyFiveItemTwo{}]~eV at \numSensCgcnnBudgetsItemThree{}, against \numSensCgcnnTableMegnetVTwoFiftyAPoint{} and \numSensCgcnnTableMegnetVTwoSixHundredFiftyFourAPoint{} [\numSensCgcnnTableMegnetVTwoSixHundredFiftyFourACiNinetyFiveItemOne{}, \numSensCgcnnTableMegnetVTwoSixHundredFiftyFourACiNinetyFiveItemTwo{}]~eV for the network on the same resample. CGCNN has \numSensCgcnnParameterCountsCgcnnCgcnnSSixHundredFiftyFour{} trainable parameters in S, fewer than the network at any tuning size (\numSensCgcnnParameterCountsMegnetVTwoSTwoHundred{} to \numSensCgcnnParameterCountsMegnetVTwoSSixHundredFiftyFour{}). With one resample the intervals reflect only the sampling of its test hosts, so the check (Figure~\ref{fig:app_sens}b) shows that the fall of $A$ with training data is not peculiar to MEGNet blocks; it sets no bound for other architectures.
